\begin{problem}{}{}{Real Mathematical Analysis, Pugh, Ex.45, p.52}
Let $(a_n)$ be a sequence of real numbers. It is \textbf{bounded} if the set $A = \{a_1, a_2, \dots \}$
is bounded. The \textbf{limit supremum}, or lim sup, of a bounded sequence $(a_n)$ as
$n \to \infty$ is
\[
    \limsup_{n \to \infty} a_n = \lim_{n \to \infty} \left( \sup_{k \ge n} a_k \right)
\]

\begin{enumerate}
    \item[(a)] Why does the $\limsup$ exist?
    \item[(b)] If $\sup \{a_n\} = \infty$, how should we define $\limsup_{n \to \infty} a_n$?
    \item[(c)] If $\lim_{n \to \infty} a_n = -\infty$, how should we define $\limsup_{n \to \infty} a_n$?
    \item[(d)] When is it true that
          \begin{align*}
              \limsup_{n \to \infty} (a_n + b_n) & \le \limsup_{n \to \infty} a_n + \limsup_{n \to \infty} b_n \\
              \limsup_{n \to \infty} c a_n       & = c \limsup_{n \to \infty} a_n ?
          \end{align*}
          When is it true they are unequal? Draw pictures that illustrate these relations.
    \item[(e)] Define the \textbf{limit infimum}, or $\liminf$, of a sequence of real numbers, and
          find a formula relating it to the limit supremum.
    \item[(f)] Prove that $\lim_{n \to \infty} a_n$ exists if and only if the sequence $(a_n)$ is bounded and
          $\liminf_{n \to \infty} a_n = \limsup_{n \to \infty} a_n$.
\end{enumerate}
\end{problem}

\begin{solution}
    \begin{itemize}
        \item[(a)]
              Define $y_n = \sup_{k \ge n} a_k$. Since $(a_n)$ is bounded,
              $\exists M \in \mathbb{R}$ such that $-M \le a_n \le M, \forall n \in \mathbb{N}$.

              Observe that $(y_n)$ is a decreasing sequence, $y_n = \sup_{k \ge n} a_k \ge \sup_{k \ge n+1}
                  a_k = y_{n+1}$ for all $n \in \mathbb{N}$.

              $(y_n)$ is bounded below by $-M$ since $y_n = \sup_{k \ge n} a_k \ge a_k \ge -M$.
              Therefore, $(y_n)$ converges and $\limsup_{n \to \infty} a_n$ exists.
        \item[(b)]
              If $\sup \{a_n\} = \infty$, then $y_n = \infty, \forall n \in \mathbb{N}$. Therefore,
              we should define $\limsup_{n \to \infty} a_n = \lim_{n \to \infty} y_n = \infty$.
        \item[(c)]
              If $\lim_{n \to \infty} a_n = -\infty$ then the sequence
              $y_n = \sup_{k \ge n} a_k$ also goes to $-\infty$. We should define
              $\limsup_{n \to \infty} a_n = -\infty$.
        \item[(d)]
              Suppose $(a_n)$ and $(b_n)$ are bounded, then it is true that $a_i + b_i
                  \le \sup_{k \ge n} a_k + \sup_{k \ge n} b_k$ for all $n \in \mathbb{N}, i \ge n$.
              This implies $\sup_{k \ge n} (a_k + b_k) \le \sup_{k \ge n} a_k + \sup_{k \ge n} b_k$.

              Therefore, $\limsup_{n \to \infty} (a_n + b_n) \le \limsup_{n \to \infty} a_n + \limsup_{n \to \infty} b_n$.

              If $c \ge 0$, $\limsup_{n \to \infty} (c a_n) = c \limsup_{n \to \infty} a_n$ because positive
              scaling preserves the order of suprema.

              If $c < 0$, $\limsup_{n \to \infty} (c a_n) = c \liminf_{n \to \infty} a_n$.
        \item[(e)]
              Suppose $(a_n)$ is a bounded sequence of real numbers. Then
              \[
                  \liminf_{n \to \infty} a_n = \lim_{n \to \infty} \left( \inf_{k \ge n} a_k \right)
              \]
              We have $\inf_{k \ge n} a_k \le \sup_{k \ge n} a_k \implies \liminf_{n \to \infty} a_n \le \limsup_{n \to \infty} a_n$.

        \item[(f)]($\Rightarrow$): Suppose $\lim_{n \to \infty} a_n$ exists, $\lim_{n \to \infty} a_n = L$. Since $(a_n)$ converges, $(a_n)$ is bounded.

              We claim that $\liminf_{n \to \infty} a_n = L$. Let $\varepsilon > 0$. $\exists N \in \mathbb{N}$ such
              that for all $k \ge N$ we have $L - \frac{\varepsilon}{2} < a_k < L + \frac{\varepsilon}{2}$.
              Since $L - \frac{\varepsilon}{2}$ is a lower bound for $\{a_k \mid k \ge N\}$,
              we have $L - \frac{\varepsilon}{2} \le \inf_{k \ge N} a_k$. For all $n \ge N$, we have:
              \[
                  L - \frac{\varepsilon}{2} \le \inf_{k \ge N} a_k \le \inf_{k \ge n} a_k \le a_n < L + \frac{\varepsilon}{2}
              \]
              Therefore, $\inf_{k \ge n} a_k \in V_\varepsilon(L)$ for all $n \ge N$. This means $\liminf_{n \to \infty} a_n = L$.
              Similarly, $\limsup_{n \to \infty} a_n = L$. We conclude that $\liminf_{n \to \infty}
                  a_n = \limsup_{n \to \infty} a_n$.

              ($\Leftarrow$): Suppose $(a_n)$ is bounded and $\limsup_{n \to \infty} a_n = \liminf_{n \to \infty} a_n = L$. 
              We will try to prove that $\lim_{n \to \infty} a_n = L$. Fix $n \in \mathbb{N}$, we have
              \[
                  \inf_{k \ge n} a_k \le a_n \le \sup_{k \ge n} a_k
              \]
              By the Squeeze Theorem, since $\lim_{n \to \infty} (\inf_{k \ge n} a_k) =
               L$ and $\lim_{n \to \infty} (\sup_{k \ge n} a_k) = L$, this implies that $\lim_{n \to \infty} a_n = L$.

    \end{itemize}
\end{solution}